\documentclass[12pt]{article}
\usepackage[utf8]{inputenc}
\usepackage[margin=1in]{geometry}
\usepackage{amsmath}
\usepackage{amsfonts}
\usepackage{amssymb}
\usepackage{enumitem}

\begin{document}

% --- PAGE 1 ---
\noindent
Lecture Scribe Generation Task \\
L15\_L17 \\
PDF PDF \\
Task: Generate a lecture scribe intended as exam-oriented reference material. \\
The scribe must be faithful to the provided context and suitable for closed-notes exam revision. \\
Context rules: Use only the attached context (lecture slides and relevant textbook excerpts). \\
Do not introduce any material, intuition, examples, or explanations that are not present in the provided context. \\
Single-call rule: Generate the complete lecture scribe in one response only. \\
Do not ask questions, clarifications, confirmations, or request additional input. Do not provide options or commentary. \\
Structural constraint: The lecture scribe must strictly follow the structure shown in the attached "Example Lecture Scribe (Structure Only)" PDF. \\
No reordering, additions, or deletions of sections are permitted. \\
The scribe must prioritize correctness over style; follow the lecture's structure and scope; include definitions, notation, assumptions, and stated results; \\
include proofs, derivations, and worked examples only if they appear in the provided context; and maintain a formal, academic tone.

\vspace{1cm}

\noindent
\textbf{Lecture Scribe: Joint Random Variables and Joint Distributions}

\section{1. Lecture Information}
\begin{itemize}[noitemsep]
    \item \textbf{Course:} CSE400 - Fundamentals of Probability in Computing.
    \item \textbf{Topic:} Lecture L15\_16\_17\_18: Joint Random Variables + Joint Distributions + Examples.
    \item \textbf{Instructor:} Dhaval Patel, PhD.
    \item \textbf{Date:} March 10, 2026.
\end{itemize}

\newpage

% --- PAGE 2 ---
\section{2. Motivation and Applications}
In science and real-world scenarios, multiple random variables are often analyzed simultaneously to understand the relationships and dependencies between them. \\
When variables are not independent, covariance and correlation are used as measures of the nature of their dependence.

\noindent
\textbf{Application Examples:}
\begin{itemize}[noitemsep]
    \item \textbf{Smart Transportation:} Analyzing Vehicle Speed (X) and Traffic Density (Y).
    \item \textbf{Wireless Network Performance:} Signal-to-Noise Ratio (SNR) (X) and Data Throughput (Y).
    \item \textbf{Weather Prediction:} Rainfall (X) and Temperature (Y).
    \item \textbf{Drone Delivery:} Wind Speed (X) and Delivery Time (Y).
\end{itemize}

\section{3. Discrete Joint Distributions}
\textbf{Setup:} Let $X$ and $Y$ be discrete random variables where $X \in \{x_1, x_2, \dots, x_n\}$ and $Y \in \{y_1, y_2, \dots, y_m\}$. \\
The ordered pair $(X, Y)$ takes values in the product set $\{(x_1, y_1), (x_1, y_2), \dots, (x_n, y_m)\}$.

\noindent
\textbf{Joint Probability Mass Function (PMF):} Defined as $p(x_i, y_j) = P(X = x_i, Y = y_j)$. \\
The Joint PMF represents the probability of the joint outcome occurring simultaneously.

\noindent
\textbf{Properties of the Joint PMF:}
\begin{itemize}
    \item The probability is bounded: $0 \le p(x_i, y_j) \le 1$
    \item The sum of all probabilities equals 1: $\sum_{i=1}^{n}\sum_{j=1}^{m} p(x_i, y_j) = 1$
\end{itemize}

\section{4. Continuous Joint Distributions}
The continuous case is analogous to the discrete case, where discrete values become continuous intervals, Joint PMFs become Joint Probability Density Functions (PDFs), and sums become double integrals.

\noindent
\textbf{Setup:} If $X \in [a, b]$ and $Y \in [c, d]$, the pair $(X, Y)$ takes values in the product region $[a, b] \times [c, d]$.

\noindent
\textbf{Joint Probability Density Function (PDF):} A function $f(x, y)$ where the probability that $(X, Y)$ lies in a small rectangle of area $dx \, dy$ is approximately $f(x, y) \, dx \, dy$.

\noindent
\textbf{Properties of the Joint PDF:}
\begin{itemize}
    \item The density is non-negative: $f(x, y) \ge 0$.
    \item The total volume under the density surface equals 1: $\int_{c}^{d} \int_{a}^{b} f(x, y) \, dx \, dy = 1$.
    \item The joint PDF $f(x, y)$ is a probability density, not a probability, meaning it may take values greater than 1.
\end{itemize}

\newpage

% --- PAGE 3 & 4 ---
\section{5. Joint Cumulative Distribution Function (CDF)}
\textbf{Definition:} For jointly distributed random variables $X$ and $Y$, the joint CDF is defined as $F(x, y) = P(X \le x, Y \le y)$.

\noindent
\textbf{Continuous Case:} The joint CDF is obtained by integrating the joint density $f(x, y)$:
\[ F(x, y) = \int_{c}^{y} \int_{a}^{x} f(u, v) \, du \, dv \]
\textbf{Recovering the Joint PDF:} Utilizing partial derivatives, the PDF can be recovered from the CDF:
\[ f(x, y) = \frac{\partial^2}{\partial x \partial y} F(x, y) \]

\noindent
\textbf{Discrete Case:} The joint CDF is obtained by summing probabilities for all valid $x_i$ and $y_j$:
\[ F(x, y) = \sum_{x_i \le x} \sum_{y_j \le y} p(x_i, y_j) \]

\noindent
\textbf{Properties of the Joint CDF:}
\begin{itemize}
    \item $F(x, y)$ is non-decreasing; it must stay constant or increase if $x$ or $y$ increase.
    \item Lower-left limit: $\lim_{(x,y) \to (-\infty, -\infty)} F(x, y) = 0$.
    \item Upper-right limit: $\lim_{(x,y) \to (\infty, \infty)} F(x, y) = 1$.
\end{itemize}

\section{6. Worked Examples}
\textbf{Example 1: Rolling Two Dice (Discrete)} \\
\textbf{Experiment:} Roll two fair dice. \\
Let $X$ be the value on the first die, and $T$ be the total sum of both dice. \\
Each specific pair of outcomes $(i, j)$ has a probability of $1/36$. \\
In the resulting joint probability table for $(X, T)$ all individual probability entries evaluate to either 0 or $1/36$. \\
The random variables $X$ and $T$ are observed to not be independent.

\vspace{0.5cm}
\noindent
\textbf{Example 2: Continuous PDF Validation and Probability Computation} \\
\textbf{Given:} $f(x, y) = 4xy$ for $0 \le x \le 1$ and $0 \le y \le 1$. The pair $(X, Y)$ jointly takes values in the unit square $[0, 1]^2$.

\noindent
\textbf{Step 1: Check Validity (Normalization)} \\
The double integral over the specified domain is:
\[ \int_{0}^{1} \int_{0}^{1} 4xy \, dx \, dy = \int_{0}^{1} \left( \int_{0}^{1} 4xy \, dx \right) dy = \int_{0}^{1} 2y \, dy = [y^2]_0^1 = 1 \]
Because $f(x, y) \ge 0$ and the total probability integrates exactly to 1, it is a valid joint PDF.

\noindent
\textbf{Step 2: Compute Probability of an Event} \\
Define the specific event $A = \{X < 0.5, Y > 0.5\}$. \\
The probability $P(A)$ is formulated as:
\[ P(A) = \int_{0}^{0.5} \int_{0.5}^{1} 4xy \, dy \, dx \]
Evaluating the integrals step-by-step:
\[ P(A) = \int_{0}^{0.5} [2xy^2]_{0.5}^{1} \, dx = \int_{0}^{0.5} 2x(1 - 0.25) \, dx = \int_{0}^{0.5} \frac{3}{2}x \, dx = \frac{3}{16} \]

\end{document}